\MakeAbstract{
本毕业设计针对电缆通道检测机器人，对比选择了轮腿构型的机器人进行研究，运动控制方面，本毕业设计首先对于轮腿进行全身建模，进行WBC（whole body control）全身控制，协同建模两条腿避免分别控制两侧腿导致倾倒。控制器设计方面，将五连杆腿简化为一条虚拟腿作为倒立摆进行控制，整体采用lqr控制器串联VMC转化为关节与轮毂电机输出。随后在matlab仿真平台上进行运动学动力学验证，机器人可连续通过四个0.25m障碍物，并保持稳定的运动状态。

在路径规划方面，本毕业设计对比了传统路径规划算法DWA与深度强化学习算法DQN和SAC，在开源环境Sparrow中进行仿真测试与参数调优。实验结果表明强化学习算法在到达率、路径效率等方面均明显优于传统路径规划算法。针对大规模数据训练加入ASL架构辅助训练，并在SAC上ji在此基础上，本毕业设计在SAC算法前串联Transformer建构，融合历史雷达数据进行训练，增强了模型应对动态障碍物的能力，TransSAC在相同训练条件下收敛速度更快，最终到达率达到0.87，路径效率也更优。经过横向对比，
}{运动控制，深度强化学习，路径规划，TransSAC}
    
\MakeAbstractEng{
This thesis investigates motion control and path planning for a wheeled-legged cable channel inspection robot operating in narrow and complex environments. First, the kinematic and dynamic models of a serial five-bar wheeled-legged robot are established and then simplified into a first-order inverted pendulum model. Based on this model, an LQR-based state-feedback controller is designed, and a VMC method is used to map the control outputs to wheel motor and joint motor inputs. MATLAB/Simulink simulations show that the proposed controller keeps the robot stable and enables it to traverse four 0.25 m obstacles successfully on both flat ground and obstacle-rich terrain.

For path planning, a reinforcement-learning simulation platform is built for dynamic obstacle scenarios, and DWA, DQN, SAC, and TransSAC are analyzed and tuned together with their reward design and distributed training framework. Simulation results indicate that the traditional DWA algorithm achieves only a 0.53 arrival rate in complex environments, with limited overall performance. After 50 million training steps, DQN stabilizes at an arrival rate of about 0.9, while TransSAC converges faster under the same training conditions and reaches an arrival rate of about 0.9 with better path efficiency, demonstrating strong navigation capability. The results provide methodological support for stable motion and intelligent navigation of cable channel inspection robots in complex environments.
}{motion control, deep reinforcement learning, path planning, TransSAC}
